% BEGIN REV09_MEMORIA_PROFUNDA
\subsection{Coordonnée profonde et récupération entière}
\label{corpus9:xii:memoria-profunda}

Le système à six coordonnées conserve le résidu exprimé et le quotient
entier qui détermine son prolongement. Nous fixons le représentant entier de
\(L_0\) présenté dans la construction régionale :
\begin{equation}
 A_H=\begin{pmatrix}
 2&2&2&1&2&1\\
 2&1&2&2&1&1\\
 1&0&0&1&0&2\\
 0&0&1&1&1&0\\
 2&2&2&0&2&1\\
 2&1&2&1&0&0
 \end{pmatrix},\qquad \det A_H=1.
 \label{corpus9:xii:matriz-profunda}
\end{equation}
Sa réduction modulo trois coïncide avec \(L_0\). Le choix du
représentant entier appartient au système de coordonnées : outre cette réduction, il fixe
l’action sur les quotients profonds. Pour des vecteurs lignes, nous définissons
\begin{equation}
 \begin{aligned}
 \mathsf{Div}_H:\mathbb Z^6&\longrightarrow
       \{0,1,2\}^6\times\mathbb Z^6,\\
 x_H&\longmapsto(u_H,x_H^+),\\
 y_H&=x_HA_H,\qquad u_H=[y_H]_3,\qquad
 x_H^+=(y_H-u_H)/3.
 \end{aligned}
 \label{corpus9:xii:division-profunda}
\end{equation}
La réduction \([\,\cdot\,]_3\) s’effectue composante par composante avec les
représentants \(0,1,2\), y compris pour les composantes entières négatives.

\begin{proposition}[Inversion exacte du résidu et du quotient
]
\label{corpus9:xii:inversion-profunda}
L’application \(\mathsf{Div}_H\) est bijective. Son inverse est
\begin{equation}
 x_H=(u_H+3x_H^+)A_H^{-1}.
 \label{corpus9:xii:inversa-profunda}
\end{equation}
Pour une trajectoire de ce système de coordonnées,
\(x_jA_H=u_j+3x_{j+1}\), la reconstruction à la profondeur
\(T\geq1\) satisfait
\begin{equation}
 x_0=\sum_{j=0}^{T-1}3^ju_jA_H^{-(j+1)}
                 +3^Tx_TA_H^{-T}.
 \label{corpus9:xii:reconstruccion-profunda}
\end{equation}
\end{proposition}
\begin{proof}
Le déterminant de \eqref{corpus9:xii:matriz-profunda} est un, de sorte que
la formule de la comatrice donne \(A_H^{-1}\in M_6(\mathbb Z)\).
La division euclidienne de chaque composante de \(x_HA_H\) détermine une
unique paire \((u_H,x_H^+)\). Réciproquement, pour tout
\((u,q)\in\{0,1,2\}^6\times\mathbb Z^6\), le vecteur
\(x=(u+3q)A_H^{-1}\) est entier et vérifie \(xA_H=u+3q\).
Son résidu est exactement \(u\) et son quotient est \(q\). Cela prouve
les deux compositions inverses. L’identité
\(x_j=(u_j+3x_{j+1})A_H^{-1}\), substituée successivement pour
\(j=0,\ldots,T-1\), donne
\eqref{corpus9:xii:reconstruccion-profunda} ; le dernier terme conserve
la coordonnée terminale qui reste après le préfixe lu.
\end{proof}

L’orientation ternaire est lue ultérieurement au moyen de
\(0\mapsto0\), \(1\mapsto+1\) et \(2\mapsto-1\), en conservant
le résidu \(u_H\) et le quotient de
\eqref{corpus9:xii:division-profunda}. Les autres champs de l’état
maintiennent leurs propres registres d’orientation, de retenue, de chemin et de bord.
La proposition~\ref{prop:ch-hensel-bruto} établit, avec ses preuves de
compatibilité et de passage à la limite, la tour arithmétique correspondant à
\(p=3\), \(d=6\) et \(A=A_H\). La réalisation des prolongements
engendrés utilise en outre les compatibilités APP--TRIT--TPK explicitées
dans la proposition~\ref{prop:cobertura-729-seccion}.

Le bloc résiduel conserve également son indice

\[
 \nu(u_H)=\sum_{j=1}^{6}(u_H)_j3^{6-j}\in\{0,\ldots,728\}.
\]
Cet indice est incorporé aux préfixes au moyen de
\(N'=729N+\nu(u_H)\). Le système de coordonnées mixte de la
section~\ref{act21:subsec:df-cilindros-mixtos} conserve simultanément
les indices ternaire et décimal et les entiers de frontière
\(A(s),B(s)\). Ses récurrences et sa compatibilité avec la troncature
sont démontrées dans la proposition~\ref{act21:prop:df-rec-frontera} et le
théorème~\ref{act21:thm:df-refinamiento-generado}. Ainsi sont liées la
récupération de la coordonnée entière et la comparaison exacte des
cylindres qui expriment ses préfixes.
% END REV09_MEMORIA_PROFUNDA
